%% paper.tex  --  LNA-Agent: News-Conditional Intraday Alpha
%% Math 285J Final Report  --  Dhruv Patel  --  June 2026

\documentclass[11pt, letterpaper]{article}

\usepackage[margin=0.95in]{geometry}
\usepackage{amsmath, amssymb, amsthm}
\usepackage{booktabs}
\usepackage{graphicx}
\usepackage[hidelinks]{hyperref}
\usepackage{natbib}
\usepackage{xcolor}
\usepackage{caption}
\usepackage{subcaption}
\usepackage{microtype}
\usepackage{multirow}
\usepackage{enumitem}
\usepackage{float}
\usepackage{parskip}
\setlength{\parskip}{4pt}

\newtheorem{definition}{Definition}

\bibliographystyle{plainnat}

\title{\textbf{News-Conditional Intraday Alpha:}\\
       LLM Sentiment, LOB Order-Flow Confirmation, and Agentic Trade Decisions\\[0.4em]
       \normalsize Math 285J Final Report}
\author{Dhruv Patel \quad {\small UCLA Mathematics} \quad
        {\small \texttt{dhruvpatel97@g.ucla.edu}}}
\date{June 2026}

\begin{document}
\maketitle

\begin{abstract}
We test whether large language model (LLM) sentiment for financial news headlines,
conditioned on limit-order-book order-flow imbalance (OFI), produces a statistically
significant Spearman information coefficient (IC) at 1-, 5-, and 15-minute forward
return horizons. Using LOBSTER level-2 data and \texttt{claude-haiku-4-5-20251001}
on a post-training-cutoff window (August--December 2025) across 10 NASDAQ equities
(2,879 matched events), the primary hypothesis does not hold: no interaction signal
survives Benjamini-Hochberg FDR control at $q = 0.05$, and the conditional OFI
regression finds no significant sentiment-direction asymmetry (Wald $p > 0.10$). The
Loughran-McDonald (LM) lexicon is the most reliable baseline, remaining nominally
significant on SPY-excess returns at 1 minute (IC $= +0.038$, $p = 0.039$).

The null, however, has structure. When we apply a single BH correction across all
subgroup tests jointly, one exploratory finding survives: on TSLA sector headlines
(EV competitor, energy, supply-chain news requiring indirect-chain reasoning), the
agent achieves IC $= +0.234$ at 15 minutes ($p = 0.024$, $n = 79$) while LM on
the same traded events scores $-0.011$. At 1 minute the gap is real but smaller:
agent $+0.284$ vs LM(traded) $+0.119$. Earnings headlines show a suggestive but underpowered signal (agent IC
$= +0.443$ at 15 min nominally, $n = 27$ traded, does not survive pooled BH). A
fair head-to-head --- scoring LM on the same traded events as the agent --- narrows
the agent vs.\ LM gap on routine news and confirms that the agent's selection
concentrates on events where LM fails. We report bootstrap CIs and pooled BH
throughout.
Code: \url{https://github.com/dhruvpatel-05/LNA-agent}
\end{abstract}

\tableofcontents
\newpage

\section{Introduction}
\label{sec:intro}

Price discovery after a public news announcement is not instantaneous. Even when a
headline is visible on a terminal, traders must read it, assess its materiality for
a specific stock, and route an order, a process that takes seconds to minutes
depending on how algorithmically saturated the name is \citep{grossklussmann2011machines,
riordan2013public}. Meanwhile the bid-ask spread temporarily widens as market makers
protect against adverse selection \citep{cont2014price}. This creates a narrow
predictive window: if sentiment direction and contemporaneous book pressure both point
the same way, their interaction should predict the next-minute return better than
either alone.

This project formalizes that idea in three layers: (1) assign a signed sentiment score
to each headline via a rule-based lexicon and an LLM; (2) measure OFI at the headline
timestamp and test whether OFI times sentiment predicts returns better than OFI or
sentiment alone; (3) pass the headline and book state to a structured-output LLM agent
that emits a trade direction, size, and horizon. The project connects directly to the
conditional-OFI framework of \citet{lu2024trade} and \citet{cont2023cross}, and to the
growing literature on LLM-based trading signals \citep{lopezlira2023chatgpt, li2025llm_kdd}.

The core pooled result is a null, but the null has structure. We contribute: (1) a
contamination-controlled evaluation pipeline that identifies where LLM inference beats
lexicon tone and where it does not; (2) a single BH correction pooled across all
subgroup tests, under which TSLA sector IC $= +0.284$ ($p < 0.001$, $n = 79$)
is the one exploratory finding that survives --- earnings is suggestive ($n = 27$,
nominally $p = 0.011$) but underpowered; (3) a fair LM-on-traded comparison showing
the agent's selection concentrates on events where LM fails hardest; and (4) LM
excess-return IC $= +0.038$ ($p = 0.039$) confirming the lexicon baseline survives
market-beta adjustment.

\section{Related Work}
\label{sec:related}

\citet{grossklussmann2011machines} show that machine-readable Reuters news implies
measurable market reactions within 1-2 minutes. \citet{riordan2013public} document
that informed order flow concentrates immediately after news disclosure in limit-order
markets. The standard OFI construction we use follows \citet{cont2014price}.
\citet{lu2024trade} extend OFI to a conditional-flow decomposition by trading-partner
type, demonstrating strong short-horizon predictability; our conditional regression is
directly analogous. \citet{cont2023cross} study cross-impact of OFI across equities,
which is relevant to our multi-ticker panel.

On the text side, \citet{loughran2011liability} build the finance-domain dictionary
we use as a point-in-time-safe baseline. \citet{tetlock2007content} establish the
media-sentiment connection to stock returns. \citet{lopezlira2023chatgpt} are among
the first to demonstrate LLM stock-return predictability and to raise the look-ahead
contamination concern formally. \citet{li2025llm_kdd} evaluate LLM investing strategies
over long horizons and find regime-dependent value. The agentic execution framework
relates to the multi-agent HFT work of \citet{mohl2025jaxmarl}. Statistical procedures
follow \citet{benjamini1995controlling} for FDR control.

\section{Data}
\label{sec:data}

\paragraph{Limit-order-book data.} We use LOBSTER level-2 reconstructions at 1-minute
bar frequency for 10 NASDAQ equities: AAPL, NVDA, AMD, META, TSLA, PEP, GILD, HON,
SOFI, and CELH. The data covers August 1 through December 31, 2025 (roughly 85 trading
days) and is licensed from \texttt{lobsterdata.com}. Each bar contains OFI computed on
the server plus 1-, 5-, and 15-minute log forward returns. SPY bars in the same format
serve as a market-return benchmark, matched at 76.6\% coverage.

\paragraph{News headlines.} Headlines are from the Alpaca news API supplemented by
a HuggingFace financial news corpus. Filtering proceeds in two steps: (1) exact-match
on ticker symbol or full company name in the headline text; (2) manual review of the
50 most-frequent match patterns per ticker to remove false positives (e.g., ``Apple''
matching fruit-market news). The final filter retains headlines where the named entity
is the primary subject, not incidental context. Timestamps are second-resolution
within the same August--December 2025 window. This curation directly affects which
events enter the IC calculation; a looser filter would dilute signal with unrelated
noise, a tighter one would reduce $n$.

\paragraph{Panel construction.} Each headline with timestamp $\tau$ is matched to the
LOB bar at $t = \lfloor \tau / 60 \rfloor$ (the bar live when the headline arrived).
Forward returns are measured from the close of that bar. After matching, the panel has
2,879 event-bar pairs. Table~\ref{tab:panel} shows composition by ticker.

\begin{table}[h]
\centering
\small
\caption{Panel composition. Low-$N$ tickers ($n < 100$) are treated as exploratory.}
\label{tab:panel}
\begin{tabular}{lrrlr}
\toprule
Ticker & Events ($n$) & Agent trades & Sector & Market cap \\
\midrule
AAPL  & 504 & 164 & Consumer electronics & Mega \\
NVDA  & 924 & 359 & Semiconductors / AI  & Mega \\
AMD   & 266 & 154 & Semiconductors       & Large \\
META  & 435 & 136 & Digital advertising  & Mega \\
TSLA  & 497 & 241 & Electric vehicles    & Large \\
PEP   &  44 &  24 & Consumer staples     & Large \\
GILD  &  53 &  27 & Biopharmaceuticals   & Large \\
HON   &  37 &  19 & Industrials          & Large \\
SOFI  &  92 &  30 & Fintech              & Small \\
CELH  &  27 &   9 & Consumer beverages   & Small \\
\midrule
Total & 2,879 & 1,163 & & \\
\bottomrule
\end{tabular}
\end{table}

\paragraph{Market context.} The evaluation window coincides with a broadly bullish
NASDAQ environment (roughly $+18\%$ over the period). This matters for the agent
evaluation in Section~\ref{sec:agent}: a systematically bearish signal will
underperform mechanically here.

\section{Methodology}
\label{sec:method}

\subsection{Sentiment Scoring}

The \textbf{Loughran-McDonald (LM)} lexicon \citep{loughran2011liability} assigns a
signed score $\mathrm{lm}_i \in \mathbb{R}$ based on net positive minus negative
finance-domain word counts, normalized by headline length. It is point-in-time safe
by construction.

The \textbf{LLM scorer} uses \texttt{claude-haiku-4-5-20251001} with temperature 0,
instructed to output a single score $\mathrm{llm}_i \in [-1,1]$ representing expected
short-term price impact. The training cutoff of this model predates August 2025,
making the evaluation window genuinely out-of-sample. The system prompt is cached
via Anthropic's \texttt{cache\_control: ephemeral} to reduce per-event cost.

\subsection{Order-Flow Imbalance}

Following \citet{cont2014price} at the best level:
\begin{equation}
    \mathrm{OFI}_t = \Delta V^B_t - \Delta V^A_t
\end{equation}
where $\Delta V^B_t$ ($\Delta V^A_t$) is the net change in best-bid (best-ask)
volume. We z-score within each trading day:
\begin{equation}
    \mathrm{ofi}_{z,t} = \frac{\mathrm{OFI}_t - \mu_d}{\sigma_d + \varepsilon},
    \qquad \varepsilon = 10^{-8}
\end{equation}
On low-volume days $\sigma_d \approx 0$ produces extreme z-scores (observed up to
$\pm 15$). We winsorize at $\pm 5\sigma$ before constructing any signal:
\begin{equation}
    \mathrm{ofi}^w_t = \mathrm{clip}(\mathrm{ofi}_{z,t},\; -5,\; +5).
\end{equation}

\subsection{Signal Definitions}

Five signals are evaluated: $s^\mathrm{ofi} = \mathrm{ofi}^w$, $s^\mathrm{lm} = \mathrm{lm}$,
$s^\mathrm{llm} = \mathrm{llm}$, and interaction terms $s^{\mathrm{ofi\times lm}} =
\mathrm{ofi}^w \cdot \mathrm{lm}$, $s^{\mathrm{ofi\times llm}} = \mathrm{ofi}^w \cdot
\mathrm{llm}$. Interaction signals are positive when OFI and sentiment agree, the core
hypothesis; continuous $\mathrm{ofi}^w$ is used rather than its sign to keep signals
on the same scale.

\subsection{IC and Bootstrap Inference}

The IC is the Spearman rank correlation of signal and forward return. We compute
percentile bootstrap CIs and a one-sided $p$-value ($B = 1{,}000$ replicates, seed 42;
see Appendix~\ref{app:stats}). Benjamini-Hochberg FDR control \citep{benjamini1995controlling}
is applied at $q = 0.05$, pooled across all 15 cells per table. A cell is BH-significant
only if it survives this correction; ``nominally significant'' means $p < 0.05$ before it.

\subsection{Conditional OFI Regression}

Following the specification of \citet{lu2024trade} with main effects included
\citep{brambor2006understanding}:
\begin{equation}
\label{eq:regression}
    r_{t+h} = \alpha + \gamma\,\mathrm{ofi}^w_t
    + \delta_1\,\mathrm{sent}^+_t + \delta_2\,\mathrm{sent}^-_t
    + \beta^+(\mathrm{ofi}^w_t \cdot \mathrm{sent}^+_t)
    + \beta^-(\mathrm{ofi}^w_t \cdot \mathrm{sent}^-_t)
    + \varepsilon_t
\end{equation}
where $\mathrm{sent}^+ = \max(\mathrm{sent}, 0)$ and $\mathrm{sent}^- = \min(\mathrm{sent}, 0)$.
Standard errors are Newey-West HAC. The primary test is a Wald test of
$H_0: \beta^+ = \beta^-$, which directly asks whether sentiment direction modulates
the OFI predictive slope. Interactions without levels would bias $\hat{\beta}^\pm$,
which is why the main effects are essential here.

\subsection{Agentic Decision Layer}

Each event is passed to a structured-output LLM agent with a ticker-specific system
prompt (sector profile plus task instructions). The agent outputs:
\[
    \mathcal{A}(h_i,\; \mathrm{lm}_i,\; \mathrm{llm}_i,\; \mathrm{ofi}^w_{t_i})
    \;\to\; \{\mathrm{direction},\; \mathrm{size} \in [0,1],\; \mathrm{horizon}\}
\]
The agent signal is $s^{\mathrm{agent}}_i = \mathrm{sign}(\mathrm{direction}) \times
\mathrm{size}$, with 0 for ``none.'' A notable design choice in the prompt: when
all signals agree, the agent should reduce conviction rather than increase it, because
consensus implies the information is already priced. This is the opposite of a naive
confirmation rule.

\section{Results}
\label{sec:results}

\subsection{Pooled IC: All Events}

Table~\ref{tab:ic_pooled} reports IC and bootstrap 95\% CIs for all five signals
pooled across 2,743--2,865 events.

\begin{table}[h]
\centering
\small
\caption{Pooled IC across all 2,879 events and all 10 tickers combined, with 95\%
bootstrap CI and Benjamini-Hochberg rejection flag ($q = 0.05$, 15 tests pooled).
$p$ = one-sided bootstrap $p$-value (fraction of replicates $\leq 0$). No cell
survives BH. $^*$ nominally significant ($p < 0.05$) before correction.}
\label{tab:ic_pooled}
\begin{tabular}{llrrrrc}
\toprule
Signal & Horizon & IC & CI$_{lo}$ & CI$_{hi}$ & $p$ & BH rej. \\
\midrule
\texttt{lm}          & 1 min  & $+0.044$ & $+0.008$ & $+0.082$ & $0.012^*$ & No \\
\texttt{lm}          & 15 min & $+0.014$ & $-0.023$ & $+0.052$ & $0.229$   & No \\
\texttt{lm}          & 5 min  & $-0.020$ & $-0.057$ & $+0.016$ & $0.862$   & No \\
\texttt{ofi}         & 5 min  & $+0.023$ & $-0.014$ & $+0.060$ & $0.110$   & No \\
\texttt{ofi}         & 1 min  & $+0.015$ & $-0.026$ & $+0.054$ & $0.234$   & No \\
\texttt{ofi}         & 15 min & $+0.014$ & $-0.027$ & $+0.054$ & $0.253$   & No \\
\texttt{ofi\_x\_llm} & 5 min  & $+0.019$ & $-0.020$ & $+0.059$ & $0.172$   & No \\
\texttt{ofi\_x\_llm} & 1 min  & $+0.003$ & $-0.036$ & $+0.044$ & $0.453$   & No \\
\texttt{ofi\_x\_llm} & 15 min & $-0.000$ & $-0.041$ & $+0.036$ & $0.517$   & No \\
\texttt{ofi\_x\_lm}  & 1 min  & $-0.032$ & $-0.065$ & $+0.003$ & $0.960$   & No \\
\texttt{ofi\_x\_lm}  & 15 min & $-0.022$ & $-0.060$ & $+0.013$ & $0.880$   & No \\
\texttt{ofi\_x\_lm}  & 5 min  & $-0.011$ & $-0.048$ & $+0.028$ & $0.710$   & No \\
\texttt{llm}         & 15 min & $-0.040$ & $-0.077$ & $+0.000$ & $0.975$   & No \\
\texttt{llm}         & 5 min  & $-0.037$ & $-0.075$ & $+0.000$ & $0.974$   & No \\
\texttt{llm}         & 1 min  & $-0.014$ & $-0.051$ & $+0.021$ & $0.783$   & No \\
\bottomrule
\multicolumn{7}{l}{\footnotesize $n = 2{,}865$ (1\,min), $2{,}836$ (5\,min), $2{,}743$ (15\,min).
$B = 1{,}000$ replicates, seed 42. $^*$nominally significant.}
\end{tabular}
\end{table}

No signal survives BH correction. LM at 1 minute is the strongest cell ($p = 0.012$),
but with 15 simultaneous tests it falls short of the BH threshold. Two structural
patterns hold across all horizons: OFI is directionally positive (range $+0.014$ to
$+0.023$) and LLM is directionally negative, consistent with the model over-reading
bearish language already priced in at arrival. Figure~\ref{fig:ic_chart} shows the
1-minute IC and bootstrap CIs for all five signals.

\textbf{Market-excess returns.} To separate directional signal from market beta in a
broadly bullish window, we subtract the matched SPY forward return from each stock
return (76.6\% of events have a SPY bar within the same minute). LM at 1 minute
remains nominally significant on excess returns (IC $= +0.038$, $p = 0.039$), nearly
unchanged from the raw IC of $+0.044$. OFI is directionally positive on excess as
well. The LLM moves from $-0.014$ (raw) to $-0.021$ (excess), a modest deterioration:
most of its negative IC is genuine directional miscalibration, not market beta.

\begin{figure}[h]
\centering
\includegraphics[width=0.82\linewidth]{fig_ic_pooled.png}
\caption{Pooled IC at 1-minute horizon for all five signals with 95\% bootstrap CIs
($n = 2{,}865$ events). LM is the only signal with a nominally positive IC;
the LLM is negative. Stars: $^* p < 0.05$, $^{**} p < 0.01$.}
\label{fig:ic_chart}
\end{figure}

\subsection{Confirmed-Event Subset}

The primary preregistered hypothesis conditions on events where LLM and OFI agree in
direction (``confirmed'' events): $\mathrm{sign}(\mathrm{llm}_i) =
\mathrm{sign}(\mathrm{ofi}_{z,i})$, both nonzero. This subset contains 979 events
(34.0\% of the panel). Table~\ref{tab:ic_confirmed} shows the result.

\begin{table}[h]
\centering
\small
\caption{IC on the confirmed-event subset: events where LLM score and OFI z-score
agree in sign (both nonzero), containing 979 events (34\% of the panel). $n$ varies
by horizon due to missing returns. No cell is significant after BH.}
\label{tab:ic_confirmed}
\begin{tabular}{llrrrr}
\toprule
Signal & Horizon & IC & CI$_{lo}$ & CI$_{hi}$ & $p$ \\
\midrule
\texttt{ofi\_x\_llm} & 5 min  & $+0.040$ & $-0.024$ & $+0.108$ & $0.115$ \\
\texttt{lm}          & 1 min  & $+0.046$ & $-0.012$ & $+0.104$ & $0.061$ \\
\texttt{ofi\_x\_llm} & 15 min & $+0.023$ & $-0.045$ & $+0.083$ & $0.248$ \\
\texttt{lm}          & 15 min & $+0.020$ & $-0.044$ & $+0.077$ & $0.286$ \\
\texttt{ofi\_x\_llm} & 1 min  & $+0.007$ & $-0.055$ & $+0.070$ & $0.430$ \\
\texttt{ofi}         & 1 min  & $-0.037$ & $-0.106$ & $+0.033$ & $0.868$ \\
\texttt{llm}         & 15 min & $-0.058$ & $-0.123$ & $+0.003$ & $0.968$ \\
\bottomrule
\end{tabular}
\end{table}

Confirmed events show \textit{weaker} OFI predictability than the full panel: OFI IC
at 5 minutes drops from $+0.023$ (full panel, $n = 2{,}836$) to $-0.037$ (confirmed
subset, $n \approx 950$). This is exactly what the calibration design anticipates:
when book pressure and LLM sentiment agree, the market has likely already acted on
the information, leaving no edge.

\subsection{Conditional OFI Regression}

Table~\ref{tab:regression} reports OLS estimates of Equation~\ref{eq:regression}
with LLM sentiment as the scorer (Newey-West HAC standard errors).

\begin{table}[h]
\centering
\small
\caption{Conditional OFI regression (Eq.~\ref{eq:regression}), pooled across all
events and tickers, using LLM score as the sentiment variable. $\hat{\gamma}$ =
OFI main-effect slope; $\hat{\beta}^+$ ($\hat{\beta}^-$) = interaction slope when
LLM score is positive (negative). Standard errors are Newey-West HAC.
Wald $p$ tests $H_0\!: \beta^+ = \beta^-$. $R^2 \approx 0.003$ at all horizons.}
\label{tab:regression}
\begin{tabular}{lrrrrrrr}
\toprule
Horizon & $n$ & $\hat{\gamma}$ & $\hat{\beta}^+$ & $\hat{\beta}^-$
        & SE$(\hat{\beta}^+)$ & SE$(\hat{\beta}^-)$ & Wald $p$ \\
\midrule
1 min  & 2,865 & $+0.000071$ & $-0.000112$ & $+0.000201$ & $0.000094$ & $0.000176$ & $0.116$ \\
5 min  & 2,836 & $+0.000028$ & $+0.000007$ & $+0.000130$ & $0.000215$ & $0.000362$ & $0.769$ \\
15 min & 2,743 & $-0.000096$ & $+0.000146$ & $+0.000299$ & $0.000413$ & $0.000809$ & $0.865$ \\
\bottomrule
\end{tabular}
\end{table}

No horizon shows a significant $\beta^+ \neq \beta^-$ ($p > 0.10$ throughout). The
interaction coefficients are numerically small (sub-basis-point per unit of OFI) and
the standard errors overlap substantially. This is the core null: OFI's predictive
content for short-horizon returns does not depend on whether the concurrent headline
is bullish or bearish, as classified by the LLM.

\subsection{Per-Ticker Heterogeneity}

Despite pooled nulls, there is meaningful cross-ticker variation. Table~\ref{tab:per_ticker}
shows selected per-ticker ICs for the two strongest baseline signals.

\begin{table}[h]
\centering
\small
\caption{Per-ticker IC at 1-minute horizon for LM lexicon and OFI (selected tickers).
IC = Spearman correlation of signal with $r_{t+1\text{min}}$, computed on all events
for that ticker. Tickers with $n < 100$ are exploratory (wide CIs expected).
$^*$ nominally significant ($p < 0.05$) before BH correction.}
\label{tab:per_ticker}
\begin{tabular}{lrrrrrr}
\toprule
Signal & Ticker & $n$ & IC & CI$_{lo}$ & CI$_{hi}$ & $p$ \\
\midrule
\multicolumn{7}{l}{\textit{LM at 1 min}} \\
\texttt{lm} & TSLA & 494 & $+0.106$ & $+0.015$ & $+0.189$ & $0.010^*$ \\
\texttt{lm} & NVDA & 918 & $+0.078$ & $+0.014$ & $+0.139$ & $0.008^*$ \\
\texttt{lm} & GILD &  53 & $+0.209$ & $-0.040$ & $+0.447$ & $0.057$   \\
\texttt{lm} & AAPL & 502 & $-0.013$ & $-0.101$ & $+0.078$ & $0.625$   \\
\texttt{lm} & AMD  & 265 & $-0.055$ & $-0.172$ & $+0.065$ & $0.831$   \\
\midrule
\multicolumn{7}{l}{\textit{OFI at 1 min}} \\
\texttt{ofi} & NVDA & 918 & $+0.058$ & $-0.011$ & $+0.136$ & $0.049^*$ \\
\texttt{ofi} & SOFI &  92 & $+0.150$ & $-0.056$ & $+0.365$ & $0.081$   \\
\texttt{ofi} & AAPL & 502 & $-0.084$ & $-0.178$ & $+0.006$ & $0.967$   \\
\bottomrule
\end{tabular}
\end{table}

NVDA (LM IC $= +0.078$, $p = 0.008$) and TSLA (LM IC $= +0.106$, $p = 0.010$) are
the two tickers where the LM lexicon is individually significant at 1 minute. NVDA
also has a nominally significant OFI IC ($+0.058$, $p = 0.049$). Both are
headline-sensitive names with moderate algorithmic coverage where price discovery takes
1-2 minutes rather than seconds. In contrast, AAPL shows LM IC $= -0.013$ and OFI
IC $= -0.084$, consistent with near-instant adjustment in a deeply liquid mega-cap.

\begin{figure}[h]
\centering
\includegraphics[width=0.82\linewidth]{fig_agent_vs_lm.png}
\caption{Agent IC vs LM IC at 1-minute horizon for all 10 tickers. Agent IC is
computed on traded events only (events where agent signal $\neq 0$, ranging from
9 to 359 per ticker); LM IC is computed on all events for that ticker. Bubble size
$\propto \sqrt{n_\text{total}}$; dark bubbles = high-$N$ tickers ($n \geq 200$).
Points below the diagonal indicate LM outperforms the agent; high-$N$ tickers
cluster below the diagonal.}
\label{fig:nvda}
\end{figure}

\section{Agentic Decision Layer}
\label{sec:agent}

\subsection{Aggregate Cross-Ticker Performance}

Table~\ref{tab:agent_cross} summarizes the agent's trading statistics and 1-minute IC
across all 10 tickers, alongside the LM baseline for comparison.

\begin{table}[h]
\centering
\small
\caption{Agentic layer performance across all 10 tickers. $n$ = total headline events
for that ticker; Trades = events where agent chose long or short (signal $\neq 0$).
Agent IC = Spearman IC of $s^{\mathrm{agent}}$ against $r_{t+1\text{min}}$ on traded
events only. Dir.\ acc.\ = fraction of trades where predicted direction matched
$\mathrm{sign}(r_{t+1\text{min}})$. LM IC = Spearman IC of LM score on all events
for that ticker at 1 minute (not just traded). Low-$N$ rows ($n < 100$) are
exploratory.}
\label{tab:agent_cross}
\begin{tabular}{lrrrrrrr}
\toprule
Ticker & $n$ & Trades & Dir.\ acc. & Long acc. & Short acc. & Agent IC & LM IC \\
\midrule
AAPL  & 504 & 164 & 45.7\% & 45.5\% & 46.3\% & $-0.037$ & $-0.013$ \\
NVDA  & 924 & 359 & 52.7\% & 54.2\% & 48.4\% & $+0.039$ & $+0.078$ \\
AMD   & 266 & 154 & 45.8\% & 50.0\% & 30.3\% & $-0.066$ & $-0.055$ \\
META  & 435 & 136 & 44.1\% & 51.1\% & 29.5\% & $-0.070$ & $-0.011$ \\
TSLA  & 497 & 241 & 49.4\% & 49.7\% & 48.9\% & $-0.012$ & $+0.106$ \\
\midrule
PEP   &  44 &  24 & 33.3\% & 29.4\% & 42.9\% & $-0.255$ & $-0.154$ \\
GILD  &  53 &  27 & 59.3\% & 60.9\% & 50.0\% & $+0.149$ & $+0.210$ \\
HON   &  37 &  19 & 42.1\% & 45.5\% & 37.5\% & $-0.028$ & $+0.072$ \\
SOFI  &  92 &  30 & 56.7\% & 62.5\% & 33.3\% & $+0.135$ & $+0.046$ \\
CELH  &  27 &   9 & 66.7\% & 75.0\% & 0.0\%  & $+0.168$ & $+0.110$ \\
\midrule
\multicolumn{8}{l}{\footnotesize High-$N$ avg ($n \geq 200$): agent IC $= -0.029$,
dir.\ acc.\ $= 47.5\%$, short acc.\ $= 40.7\%$.} \\
\multicolumn{8}{l}{\footnotesize All-ticker avg: agent IC $= +0.002$,
dir.\ acc.\ $= 49.6\%$, short acc.\ $= 36.7\%$.} \\
\bottomrule
\end{tabular}
\end{table}

Among the five high-$N$ tickers, only NVDA has positive agent IC ($+0.039$). The
high-$N$ average is $-0.029$, and on TSLA the LM baseline substantially outperforms
the agent in aggregate ($+0.106$ vs $-0.012$). The pattern driving the losses is
short-direction bias: pooled short accuracy is 36.7\%, well below the 50\% baseline,
while long accuracy is near 50\% for most tickers. AMD and META are the worst cases
(short accuracy 30.3\% and 29.5\%). The LLM was trained on general internet text
where cautionary language is common; in a broadly bullish market this produces a
systematic bearish miscalibration. The aggregate agent IC is therefore a convolution of two things:
headline classification skill and directional calibration bias. Separating them requires
stratifying by event type, which Section~\ref{sec:eventtype} does.
Figure~\ref{fig:nvda} shows where each ticker sits on the agent-vs-LM IC scatter at
1 minute; the high-$N$ tickers (dark bubbles) cluster below the diagonal, confirming
that LM outperforms the agent in aggregate for the names where statistical power is
highest.

\subsection{Task-Type Stratification: Where Each Method Wins}
\label{sec:eventtype}

We assign each headline a keyword-based event type:
\textit{earnings} (quarterly results), \textit{analyst} (upgrades/downgrades/targets),
\textit{M\&A} (deals, acquisitions), \textit{legal} (regulatory, lawsuits), or
\textit{other} (the residual 83\% of events). Table~\ref{tab:eventtype} reports agent
and LM IC at 1 and 15 minutes for each type, with two key methodological improvements:
(1) a \textbf{single BH correction pooled across all cells in this table and
Table~\ref{tab:relevance}} (no cell in this table survives it); and (2) \textbf{LM
scored on the same traded subset as the agent} (\textit{LM traded}), enabling a fair
head-to-head rather than scoring LM on all events and agent on traded only.

\begin{table}[h]
\centering
\small
\caption{Agent IC and LM IC stratified by keyword-classified event type, pooled across
all 10 tickers. $n_{\text{all}}$ = all events of this type; $n_{\text{trd}}$ = events
where the agent traded (signal $\neq 0$). Agent IC is on traded events; LM(all) is
on all events (as in prior work); LM(traded) is on the same traded subset for a fair
comparison. \textbf{No cell in this table survives} the single pooled BH correction
applied jointly with Table~\ref{tab:relevance}; all values shown are nominal.
$^{\cdot}$ nominal $p < 0.10$; no stars are shown to avoid confusion with BH rejection.}
\label{tab:eventtype}
\begin{tabular}{lrrrrrrrr}
\toprule
Type & $n_{\text{all}}$ & $n_{\text{trd}}$ & Agent 1m & Agent 15m & LM(all) 1m & LM(trd) 1m & LM(trd) 15m & Winner \\
\midrule
earnings &  82 & 27 & $+0.288^{\,\cdot}$ & $+0.443^{\,\cdot}$ & $-0.207$ & $-0.310$ & $-0.136$ & Agent \\
legal    &  36 & 23 & $+0.160$            & ---                 & $-0.039$ & $-0.093$ & ---      & Agent \\
analyst  & 237 &192 & $+0.020$            & $+0.001$            & $+0.065$ & $+0.005$ & $+0.003$ & Agent \\
M\&A     & 103 & 56 & $-0.064$            & $-0.167$            & $+0.084$ & $-0.034$ & $+0.331^{\,\cdot}$ & LM \\
other    &2,400&852 & $-0.047$            & $-0.003$            & $+0.056^{\,\cdot}$ & $+0.034$ & $+0.010$ & LM  \\
\midrule
\multicolumn{9}{l}{\footnotesize Legal 15m and analyst-15m LM(trd) excluded ($n < 30$). Guidance ($n=21$) excluded throughout.} \\
\bottomrule
\end{tabular}
\end{table}

The task-type direction is consistent: on events requiring numerical or structural
inference (earnings, legal), the agent IC is positive and LM is negative. On events
where tone is the signal (M\&A, other), LM leads.

The fair comparison (LM traded vs.\ agent) changes the picture meaningfully. For
analyst events, LM(all) IC $= +0.065$ suggests LM wins; but LM(traded) $= +0.005$
shows that advantage nearly disappears when LM is restricted to the agent's traded
events. Similarly, for \textit{other} events, LM(all) $= +0.056$ drops to LM(traded)
$= +0.034$. The agent's abstention has signal value: it selects events where LM is
weakest. For earnings, LM(traded) $= -0.310$ --- more negative than LM(all) $= -0.207$
--- confirming the agent focuses on events where LM fails most.

The earnings 15-min cell ($+0.443$, $p = 0.011$ nominally, $n = 27$) is suggestive
but underpowered. At $n = 27$, the Spearman SE $\approx 0.20$; the 1-min cell already
has $p = 0.072$. The same events are evaluated at both horizons, so these are not
independent tests. Leave-one-out shows the direction is consistent
($+0.408$--$+0.530$ at 15 min across ticker holdouts), but these resample the same
small pool and test ticker concentration, not statistical power. We report the earnings
finding as exploratory.

\subsection{Indirect Reasoning: The TSLA Sector Case}
\label{sec:relevance}

The agent classifies each headline as \textit{direct} (company or ticker explicitly
named), \textit{sector} (competitor, supplier, or industry news with an indirect
link), \textit{macro} (broad market or rates with a more diffuse link), or
\textit{unrelated}. This classification is itself an output of the agent's reasoning.
Stratifying IC by class reveals a clean split, shown in Table~\ref{tab:relevance}.

\begin{table}[h]
\centering
\small
\caption{Agent IC and LM IC (on the same traded events) stratified by the agent's
own relevance classification. $n$ = agent trades (signal $\neq 0$) in that cell.
POOLED = all 10 tickers combined. A single BH correction is applied jointly across
this table and Table~\ref{tab:eventtype}: \textbf{TSLA sector at 1 min is the only
cell that survives} ($\dagger$). All other $p$-values are nominal.
$^* p < 0.05$, $^{**} p < 0.001$; $\dagger$ survives pooled BH ($q = 0.05$).}
\label{tab:relevance}
\begin{tabular}{llrrrrrrr}
\toprule
Rel. & Ticker & H & $n$ & Agent IC & CI$_{lo}$ & CI$_{hi}$ & $p$ & LM(trd) IC \\
\midrule
\multicolumn{9}{l}{\textit{Sector headlines}} \\
sector & TSLA   & 1 min  & 79 & $+0.284^{\dagger}$ & $+0.102$ & $+0.464$ & $<0.001$ & $+0.119$ \\
sector & TSLA   & 15 min & 79 & $+0.234^*$         & $+0.008$ & $+0.445$ & $0.024$  & $-0.011$ \\
sector & POOLED & 1 min  &476 & $+0.006$            & $-0.088$ & $+0.086$ & $0.470$  & $+0.020$ \\
\midrule
\multicolumn{9}{l}{\textit{Direct headlines}} \\
direct & NVDA   & 1 min  &154 & $+0.125^*$ & $-0.015$ & $+0.270$ & $0.049$ & $+0.068$ \\
direct & POOLED & 1 min  &655 & $-0.006$   & $-0.080$ & $+0.069$ & $0.570$ & $-0.009$ \\
\bottomrule
\end{tabular}
\end{table}

\textbf{TSLA sector is the headline finding} and the only cell that survives the full
pooled BH correction. Agent IC $= +0.284$ at 1 minute ($p < 0.001$, $n = 79$) and
$+0.234$ at 15 minutes ($p = 0.024$, $n = 79$). LM on the same traded events scores
$+0.119$ at 1 minute (nominally positive but not BH-significant) and $-0.011$ at 15
minutes --- the agent's advantage grows with horizon, consistent with slower-moving
indirect information.

The mechanism is the same as the earnings case: TSLA's price responds to EV sector
dynamics, competitor moves, and energy-policy headlines. An LM lexicon cannot reason
through the chain from ``competitor cuts prices'' to ``TSLA margin pressure''; the
agent can. The pooled sector IC across all tickers ($+0.006$, $n = 476$) is near zero,
confirming this is TSLA-specific: not every stock has sector dynamics this central to
its price. Surviving pooled BH is necessary but not sufficient evidence --- this is
one post-hoc cell out of roughly 25 tested at $q = 0.05$, the subset is defined by
the agent's own relevance label (so labeling and trading are coupled), and the effect
has not been replicated out-of-sample. The real test is whether other narrative-driven
names (EV peers, heavily-shorted names with dense indirect-news flow) show the same
pattern. NVDA direct headlines ($n = 154$, agent IC $= +0.125$, $p = 0.049$
nominal) are a candidate but do not survive BH.
Figure~\ref{fig:tsla_relevance} shows both horizons.

\begin{figure}[h]
\centering
\includegraphics[width=\linewidth]{fig_tsla_relevance.png}
\caption{\textit{Left}: Agent IC for TSLA headlines classified as \textit{sector} (EV
competitor, energy, supplier news with indirect TSLA link) vs \textit{direct} (TSLA
explicitly named) at 1-minute and 15-minute horizons. Error bars are 95\% bootstrap
CIs; $n = 79$ TSLA sector trades, $n = 156$ TSLA direct trades.
\textit{Right}: The two strongest individual positive agent cells in the study ---
TSLA sector at 1-min (IC $= +0.284$, $p < 0.001$, $n = 79$, survives pooled BH)
and NVDA direct at 1-min (IC $= +0.125$, $p = 0.049$, $n = 154$) --- with
$p$-values annotated. Both are agent-traded events only.}
\label{fig:tsla_relevance}
\end{figure}

\subsection{Signal Agreement and the Calibration Design}

A second diagnostic uses the agent's own \texttt{signal\_agreement} label, which
records whether OFI, LM, and LLM were all pointing the same direction (``confirms''),
in disagreement (``conflicts''), or mixed (``neutral'') at the time of the trade.
The agent prompt was designed with an unusual calibration: when all signals agree,
reduce conviction, because consensus implies the information is already priced. The
results in Table~\ref{tab:agreement} validate this design.

\begin{table}[h]
\centering
\small
\caption{Agent IC on traded events, stratified by the agent's \texttt{signal\_agreement}
label at the time of the trade. \textit{confirms} = OFI, LM, and LLM all point the
same direction as the agent's chosen trade; \textit{conflicts} = at least one
disagrees; \textit{neutral} = mixed. $n$ = number of agent trades in that group.
IC = Spearman correlation of agent signal with $r_{t+1\text{min}}$, pooled across
all tickers.}
\label{tab:agreement}
\begin{tabular}{lrrrrr}
\toprule
Agreement & $n$ & IC & CI$_{lo}$ & CI$_{hi}$ & $p$ \\
\midrule
conflicts & 1,086 & $-0.035$ & $-0.097$ & $+0.026$ & $0.866$ \\
neutral   &    11 & $+0.020$ & \multicolumn{2}{c}{(underpowered)} & $0.472$ \\
confirms  &    42 & $-0.265$ & $-0.545$ & $+0.063$ & $0.949$ \\
\bottomrule
\end{tabular}
\end{table}

The ``confirms'' group ($n = 42$) has the worst IC ($-0.265$): when OFI, LM, and LLM
all agreed, the market had already acted and any remaining trade was adversely selected.
The agent suppressed most high-agreement trades by design, but the few that slipped
through confirm the calibration logic was correct.

\section{Discussion}
\label{sec:discussion}

\paragraph{The primary null.} The conditional regression finds no significant
$\beta^+ \neq \beta^-$, meaning news sentiment does not measurably modulate the OFI
predictive slope. The confirmed-event IC is weaker than the full-panel IC for OFI,
not stronger. These are proper nulls: the pooled sample of 2,865 events can detect an
IC around $\pm 0.037$ at 80\% power, so the null is not a power failure for the
primary test. The negative pooled LLM IC ($-0.014$ at 1 minute) also argues against
look-ahead contamination: a model that memorized outcomes would show spuriously
\emph{positive} IC, not negative. The evaluation window (August--December 2025)
additionally postdates the \texttt{claude-haiku-4-5} training cutoff, making these
events genuinely out-of-sample \citep{lopezlira2023chatgpt}.

\paragraph{Task-type structure and statistical limits.} Section~\ref{sec:agent}
reveals a consistent directional pattern: the agent leads on events requiring
numerical or structural inference (earnings, sector); LM leads on events where tone
is the signal (analyst, M\&A, other). A fair comparison --- LM scored on the same
traded subset as the agent --- narrows the LM advantage on routine events considerably
(LM(traded) $= +0.034$ on \textit{other} events vs.\ LM(all) $= +0.056$), and shows
the agent's abstention concentrates on events where LM fails hardest. When a single
BH correction is applied across all subgroup tests jointly, only TSLA sector at 1
minute survives ($+0.284$, $p < 0.001$). The earnings 15-min cell is suggestive
($+0.443$, $p = 0.011$ nominal, $n = 27$) but at $n = 27$ the Spearman SE is $\approx
0.20$ and both horizons test the same events, so we characterize it as exploratory.

\paragraph{The short-direction bias.} The aggregate agent underperformance is partly
attributable to a calibration error: pooled short accuracy is 36.7\%, well below the
50\% baseline. Market-excess return analysis (subtracting matched SPY returns) shows
the LLM IC moves from $-0.014$ (raw) to $-0.021$ (excess) at 1 minute --- most of
the negative IC is genuine directional miscalibration, not market beta. The LLM
over-reads cautionary financial language across all conditions. LM excess IC remains
nominally significant ($+0.038$, $p = 0.039$), confirming that the lexicon baseline
is not a beta artefact.

\paragraph{Relationship to prior work.} \citet{lu2024trade} demonstrate strong
short-horizon predictability from conditionally signed OFI; the null interaction here
suggests OFI's predictive content is captured by book dynamics alone and does not
depend on concurrent news direction. The 1-2 minute price discovery window we observe
for NVDA and TSLA is consistent with \citet{grossklussmann2011machines}'s finding that
headline-driven names that are not fully algorithmically saturated show slower adjustment.

\section{Paper PnL}
\label{sec:pnl}

Table~\ref{tab:pnl} reports gross paper PnL pooled across all events.
All results are pre-cost; mean-per-trade is sub-basis-point for every signal.

\begin{table}[h]
\centering
\small
\caption{Gross paper PnL for a sign-of-signal strategy applied pooled across all 10
tickers and all matched events. Strategy: go long (position $= +1$) if signal $> 0$,
short (position $= -1$) if signal $< 0$, flat if signal $= 0$. Total PnL $=$ sum
of position $\times$ forward log-return across all $n$ events. No transaction costs
applied. Mean/trade and Total PnL are in log-return units. Hit rate = fraction of
non-flat trades with correct direction sign.}
\label{tab:pnl}
\begin{tabular}{llrrrr}
\toprule
Signal & Horizon & Total PnL & Mean/trade & Hit rate & $n$ \\
\midrule
\texttt{ofi}         & 5 min  & $+0.257$ & $+0.000090$ & 50.4\% & 2,836 \\
\texttt{ofi}         & 15 min & $+0.141$ & $+0.000052$ & 50.6\% & 2,743 \\
\texttt{ofi\_x\_llm} & 5 min  & $+0.105$ & $+0.000037$ & 49.9\% & 2,836 \\
\texttt{lm}          & 1 min  & $+0.050$ & $+0.000017$ & 54.6\% & 2,865 \\
\texttt{lm}          & 15 min & $+0.044$ & $+0.000016$ & 51.6\% & 2,743 \\
\texttt{ofi}         & 1 min  & $+0.029$ & $+0.000010$ & 50.1\% & 2,865 \\
\midrule
\texttt{llm}         & 5 min  & $-0.248$ & $-0.000087$ & 47.9\% & 2,836 \\
\texttt{llm}         & 15 min & $-0.385$ & $-0.000140$ & 48.5\% & 2,743 \\
\bottomrule
\end{tabular}
\end{table}

OFI at 5 minutes leads ($+0.257$ total, 0.9 bp/trade), but bid-ask spreads widen
on news arrival, so net-of-cost performance is likely zero or negative.

\section{Limitations and Future Work}
\label{sec:limits}

Five months of data is the binding constraint: minimum detectable IC is $\pm 0.09$
at $n = 500$ and $\pm 0.20$ at $n = 100$, so many per-ticker null results may be
power failures rather than true zeros. All PnL results are gross of costs; the
mean-per-trade is sub-basis-point and bid-ask spreads widen on news arrival, making
net-of-cost performance likely zero or negative --- a net-of-cost backtest using
the LOBSTER book snapshot at each event timestamp is the single most important missing
piece before economic interpretation. OFI is computed at the best level only, despite
L2 data being available; deeper-book OFI typically carries more signal
\citep{cont2014price} and is a natural extension. Natural further extensions include
cross-ticker spillover and longer post-cutoff evaluation windows to improve power on
per-type strata.

\section{Conclusion}
\label{sec:conclusion}

We test a three-layer news-plus-LOB architecture on 2,879 matched events across 10
NASDAQ equities in a post-training-cutoff evaluation window. The primary hypothesis
--- that news sentiment modulates the OFI predictive slope --- is not supported. The
pooled result is a null, but the null has a consistent internal structure that is
itself the main finding.

Under a single BH correction pooled across all subgroup tests, TSLA sector IC
$= +0.284$ at 1 minute ($p < 0.001$, $n = 79$) is the one exploratory finding that
survives. The agent selects sector headlines where indirect chain-reasoning is required
--- and achieves consistent positive IC across both horizons (agent $+0.284$ vs
LM(traded) $+0.119$ at 1 min; agent $+0.234$ vs LM(traded) $-0.011$ at 15 min). Earnings shows a suggestive direction (agent IC
$= +0.443$ at 15 min nominally) but $n = 27$ is underpowered at Spearman SE $\approx
0.20$; we report it as exploratory.

A fair comparison --- scoring LM on the same traded events as the agent --- changes
the picture on routine news: LM(traded) IC on \textit{other} events is $+0.034$, down
from LM(all) $= +0.056$. The agent's abstention has signal value independent of its
direction decisions. LM excess IC $= +0.038$ ($p = 0.039$) confirms the lexicon
baseline survives market-beta adjustment. The LLM's short-direction bias
(36.7\% short accuracy) is mostly genuine miscalibration, not beta artefact (LLM
excess IC $= -0.021$, worse than raw $-0.014$).

The practical takeaway is a task-type taxonomy, not an alpha claim. Use LM where
word tone is informative; route to an LLM agent on events requiring numerical or
indirect-chain reasoning. The net-of-cost hurdle --- bid-ask spreads widen on news
arrival and all mean-per-trade PnL is sub-basis-point gross --- is the outstanding
empirical question, and the single most important extension before economic
interpretation of these results.

\section*{Code and Reproducibility}

Full pipeline code: \url{https://github.com/dhruvpatel-05/LNA-agent}. The pipeline runs via
\texttt{src/run.py} with subcommands \texttt{sentiment}, \texttt{agent}, and
\texttt{analyze}. Random seeds are fixed at 42 for all bootstrap procedures. The
model version (\texttt{claude-haiku-4-5-20251001}) and temperature (0) are pinned
in \texttt{src/config.py}.

\appendix

\section{LLM Prompts}
\label{app:prompt}

\paragraph{Sentiment scorer.}
\begin{verbatim}
System: You are a financial news sentiment analyst.
Output a single decimal number in [-1.0, 1.0].
No refusals, no explanations, no other text -- one number only.
Score the expected short-term stock price impact:
-1.0 = clearly bearish, 0.0 = neutral, +1.0 = clearly bullish.

User: [HEADLINE]
\end{verbatim}

\paragraph{Agent calibration note (excerpt from system prompt).}
\begin{verbatim}
CALIBRATION:
When signals confirm each other (OFI, LM, LLM all agree), the
information is likely already reflected in prices. High agreement
should make you more cautious -- reduce size.
When signals conflict, you may be resolving genuine ambiguity. Clear
headline + direct company link + conflicting signals warrants higher
conviction.
\end{verbatim}

Full agent prompt is in \texttt{src/config.py:AGENT\_TASK}.

\section{OFI Construction}
\label{app:ofi}

For bar $t$, let $(P^B_t, V^B_t)$ and $(P^A_t, V^A_t)$ denote best-bid and best-ask
price/volume at the bar close. Define:
\begin{align}
    e^B_t &= V^B_t \cdot \mathbf{1}[P^B_t \geq P^B_{t-1}]
             - V^B_{t-1} \cdot \mathbf{1}[P^B_t \leq P^B_{t-1}] \\
    e^A_t &= V^A_t \cdot \mathbf{1}[P^A_t \leq P^A_{t-1}]
             - V^A_{t-1} \cdot \mathbf{1}[P^A_t \geq P^A_{t-1}]
\end{align}
Then $\mathrm{OFI}_t = e^B_t - e^A_t$ \citep{cont2014price}. Winsorization at
$\pm 5\sigma_d$ removes artifacts from degenerate days.

\section{Bootstrap IC and BH Procedure}
\label{app:stats}

\paragraph{Bootstrap IC.} For $n$ pairs $\{(s_i, r_i)\}$: compute
$\widehat{\mathrm{IC}} = \rho_S(s,r)$; draw $B = 1{,}000$ bootstrap samples
(seed 42); compute $\mathrm{IC}^{(b)}$ for each; report percentile CI and
$p = B^{-1}\sum_b \mathbf{1}[\mathrm{IC}^{(b)} \leq 0]$.

\paragraph{Benjamini-Hochberg.} Given $m$ tests with sorted $p$-values
$p_{(1)} \leq \cdots \leq p_{(m)}$, reject $H_{(k)}$ for all $k \leq k^*$ where
$k^* = \max\{k: p_{(k)} \leq kq/m\}$. We apply BH once per table with $m = 15$.

\bibliography{refs}

\end{document}
